\documentclass{article}
\usepackage[utf8]{inputenc}
\usepackage[T1]{fontenc}
\usepackage{geometry}
\geometry{
    a4paper,
    left=2.5cm,
    right=2.5cm,
    top=2.5cm,
    bottom=2.5cm
}
\usepackage{graphicx}
\usepackage{float}
\usepackage{hyperref}
\usepackage[english]{babel}
\usepackage{amsmath}
\usepackage{amssymb}
\usepackage{array}
\usepackage{booktabs}
\usepackage{caption}
\usepackage{subcaption}
\usepackage{enumitem}
\usepackage{fancyhdr}
\usepackage{titlesec}
\usepackage{url}
\usepackage{textcomp}
\usepackage{xcolor}
\usepackage{listings}
\usepackage{multirow}
\usepackage{longtable}
\usepackage{tabularx}
\usepackage{times}

% Set up page style
\pagestyle{fancy}
\fancyhf{}
\rfoot{\thepage}
\renewcommand{\headrulewidth}{0pt}

% Title formatting
\titleformat{\chapter}[display]{\normalfont\bfseries}{\chaptertitle\thechapter}{0pt}{\Huge}
\titlespacing*{\chapter}{0pt}{50pt}{40pt}

% Document info
\title{Template Document}
\author{Academic Support Agent}
\date{\today}

\begin{document}
\maketitle

%----------- Informations du rapport ---------

\titre{Titre du rapport} %Titre du fichier .pdf
\UE{UE PRO} %Nom de la UE
\sujet{\LaTeX Approfondi} %Nom du sujet

\enseignant{Pr nom \textsc{Nom}} %Nom de l'enseignant

\eleves{Pr nom \textsc{Nom} \\
		Pr nom \textsc{Nom} \\ 
		Pr nom \textsc{Nom} } %Nom des  l ves

%----------- Initialisation -------------------
        
\fairemarges %Afficher les marges
\fairepagedegarde %Cr er la page de garde
\tabledematieres %Cr er la table de mati res

%------------ Corps du rapport ----------------


\section{Premi re section} 

\lipsum[3-4]%Effacer cette ligne et  crire le texte souhait 

\subsection{Subsection}

\lipsum[3-4] %Effacer cette ligne et  crire le texte souhait 

\section{Deuxi me section}

\lipsum[3-5] %Effacer cette ligne et  crire le texte souhait 

%------------- Commandes utiles ----------------

\section{Quelques commandes}

Voici quelques commandes utiles :

%------ Pour ins rer et citer une image centralis e -----

\insererfigure{logos/Lyon 3 logo.png}{3cm}{L gende de la figure}{Label de la figure}
% Le premier argument est le chemin pour la photo
% Le deuxi me est la hauteur de la photo
% Le troisi me la l gende
% Le quatri me le label
Ici, je cite l'image \ref{fig: Label de la figure}


%------- Pour ins rer et citer une  quation --------------

\begin{equation} \label{eq: exemple}
\rho + \Delta = 42
\end{equation}

L' quation \ref{eq: exemple} est cit  ici. 

% ------- Pour  crire des variables ----------------------

Pour  crire des variables dans le texte, il suffit de mettre le symbole \$ entre le texte souhait  comme : constante $\rho$.
\end{document}